% Vietnamese translation for OI Public Library
% Source-PDF: IOI中国国家候选队论文/2008/Day1/7.方戈《浅析信息学竞赛中一类与物理有关的问题》/ars longa.pdf
\documentclass[11pt]{article}
\usepackage{oipl-vn}

\title{Ars Longa}
\author{World Finals -- San Antonio -- Texas -- 2005/2006}
\date{}

\begin{document}
\maketitle

\origtitle{3563 -- Ars Longa}
\sourcepath{IOI-China-Candidate-Team-Papers/2008/Day1/7-Fang-Ge/ars longa.pdf}

\section*{Đề bài}

Bạn vừa được cảm hứng đánh trúng và đang thiết kế một tác phẩm điêu khắc nghệ
thuật mới thật đẹp cho sảnh của bảo tàng địa phương. Vì những lý do nghệ thuật
rất quan trọng, bạn muốn dùng đúng một số vật liệu rất cụ thể. Tuy nhiên, bạn
không chắc vật lý có đứng về phía mình hay không. Liệu tác phẩm điêu khắc của
bạn có thật sự đứng vững được không?

Tác phẩm sẽ gồm nhiều khớp cầu, mỗi khớp nặng $1$ kilogram, và nhiều thanh nối
các khớp ấy. Các thanh có khối lượng không đáng kể. Thanh không thể bị kéo giãn
hay nén ngắn, và cũng không bao giờ rời khỏi khớp. Tuy vậy, chúng có thể tự do
quay quanh khớp theo mọi hướng. Các khớp nằm trên mặt đất được dán cố định tại
chỗ; mọi khớp khác đều có thể chuyển động. Để đơn giản, bạn có thể bỏ qua ảnh
hưởng của việc các thanh giao nhau; mỗi thanh chỉ tác dụng lực lên đúng $2$
khớp mà nó nối. Ngoài ra, mỗi khớp ở trên không đều có ít nhất một thanh đi ra
không song song với mặt đất. Điều này loại bỏ trường hợp suy biến khi một quả
cầu chỉ được một cấu trúc cứng chống đỡ theo phương ngang. Trong đời thực, nó
sẽ võng xuống một chút.

Hãy viết chương trình xác định cấu trúc của bạn có \term{tĩnh} hay không,
nghĩa là nó có lập tức chuyển động dưới tác dụng của trọng lực hay không. Lưu
ý rằng mỗi thanh có thể truyền một lực căng lớn tùy ý dọc theo chiều dài của
nó, và ``tĩnh'' nghĩa là tại mỗi khớp, các lực căng cân bằng đúng trọng lượng
của khớp đó.

Nếu cấu trúc là tĩnh, bạn còn phải xác định nó có \term{ổn định} hay không,
nghĩa là nó có chuyển động khi bị nhiễu loạn nhẹ bằng cách kéo các khớp của nó
hay không.

\section*{Input}

Dữ liệu vào chứa nhiều mô tả tác phẩm điêu khắc. Mỗi mô tả bắt đầu bằng một
dòng chứa hai số nguyên $J$ và $R$, lần lượt là số khớp và số thanh trong cấu
trúc. Các khớp được đánh số từ $1$ đến $J$. Tiếp theo là $J$ dòng, mỗi dòng
ứng với một khớp và chứa $3$ số thực cho biết tọa độ $x$, $y$, $z$ của khớp
đó. Sau đó là $R$ dòng, mỗi dòng chứa $2$ số nguyên phân biệt cho biết hai
khớp được nối bởi một thanh.

Mỗi thanh có đúng độ dài cần thiết để nối hai khớp của nó. Tọa độ $z$ luôn
không âm; tọa độ $z$ bằng $0$ cho biết khớp nằm trên mặt đất và được cố định
tại chỗ. Có không quá $100$ khớp và $100$ thanh.

Sau mô tả cuối cùng là một dòng chứa hai số không.

\section*{Output}

Với mỗi mô tả tác phẩm điêu khắc, in ra \texttt{NON-STATIC},
\texttt{UNSTABLE}, hoặc \texttt{STABLE}, như trong output mẫu.

\section*{Sample Input}

\begin{verbatim}
4 3
0 0 0
-1.0 -0.5 1.0
1.0 -0.5 1.0
0 1.0 1.0
1 2
1 3
1 4
4 6
0 0 0
-1.0 -0.5 1.0
1.0 -0.5 1.0
0 1.0 1.0
1 2
1 3
1 4
2 3
2 4
3 4
7 9
0 0 0
-1.0 -0.5 1.0
1.0 -0.5 1.0
0 1.0 1.0
-1.0 -0.5 0
1.0 0.5 0
0 1.0 0
1 2
1 3
1 4
2 3
2 4
3 4
2 5
3 6
4 7
0 0
\end{verbatim}

\section*{Sample Output}

\begin{verbatim}
Sculpture 1: NON-STATIC
Sculpture 2: UNSTABLE
Sculpture 3: STABLE
\end{verbatim}

\section*{Ghi chú nguồn}

Bộ dữ liệu dùng trong bài này là dữ liệu không chính thức do Derek Kisman chuẩn
bị. Vì vậy, mọi sai sót ở đây không hàm ý sai sót trong dữ liệu của hệ chấm
chính thức. Chỉ Derek Kisman chịu trách nhiệm về các sai sót. Hãy báo lỗi đến
\texttt{dkisman@acm.org}.

Người tạo bộ test: Derek Kisman.

\end{document}
